\chapter{Discussion and Conclusion}
\label{chap:discussion-conclusion}

This chapter interprets the findings relative to the research questions and objectives. The dissertation evaluated the performance of
Supervised Machine Learning Models (Random Forest, XGBoost, DNN) and Unsupervised Machine Learning Models (Isolation Forest, One-Class SVM, and improved 
Auto Encoder) on CICIDS2017, CSE-CIC-IDS2018, and UNSW-NB15 for cross-dataset generalization for NIDS.

Two aspects of performance must be distinguished. ROC-AUC measures how well continuous anomaly scores rank attacks above benign flows 
across thresholds. F1\_score and False-Positive Rate (FPR) describe the detector's behaviour at the threshold selected using source 
validation data. A model can retain useful score discrimination after transfer while producing too many false alerts at that frozen 
thresholds. The implication, contribution, limitations, and conclusions are considered with this distinction in mind.

\section{Implications}
\label{sec:conclusion-implications}

\subsection*{Implications for Cross-Dataset Generalization}

The within-dataset results provided useful baselines but did not predict reliable performance on external datasets. We performed within dataset
generalization using Supervised Models only for CICIDS2017. All the models generated F1\_Score and ROC-AUC higher than 0.99; XGBoost being the best 
model. Isolation Forest achieved the highest within-dataset F1\_Score on CICIDS2017 (0.7250) and CSE-CIC-IDS2018 (0.5700), while the improved 
autoencoder led on UNSW-NB15 (0.8112). Meanwhile within-dataset ROC-AUC was 0.9097 for Isolation Forest, 0.7914 for One-Class SVM, and 0.8335 for the 
autoencoder. The corresponding cross-dataset means were 0.6250, 0.6211, and 0.6490. These are arithmetic averages across different paths, not 
estimates of performance variation across repeated runs.  All the models (Supervised and Unsupervised) had lower mean ROC-AUC and F1\_Score 
across the cross-dataset paths than across their within-dataset paths.

They nevertheless describe a consistent deterioration in the recorded experiment. The results reinforces the need for external testing highlighted 
by \cite{Verkerken2022} and the broader cross-dataset concerns discussed by \cite{Cantone2024}. Transfer direction also mattered. A model trained 
on CICIDS2017 and evaluated on CSE-CIC-IDS2018 did not necessarily behave like a model trained on CSE-CIC-IDS2018 and evaluated on CICIDS2017. 
Reversing the direction changes the normal training distribution, fitted preprocessing parameters, selected configuration, and source-calibrated 
threshold. Cross-dataset findings should therefore name both the source and target rather than treat a dataset pair as one undirected experiment.

\subsection*{Implications for Intrusion Detection Practice}

The findings from the comparative analysis illustrate that 'Accuracy' or 'ROC-AUC' doesn't attribute as the effective indicators of the perfomance of 
the models. Although some models showed significant ROC-AUC scores, their operational thresholds and performance were limited in most cases. For example, 
Isolation Forest trained on UNSW-NB15 and tested on CICIDS2017 yielded a ROC-AUC score of 0.8029, yet the threshold calibration gave a false positive rate 
of 0.9307, meaning that it wrongly flagged 59786 out of 64240 legitimate flows. Similarly, the Improved Autoencoder trained on CSE-CIC-IDS2018 yielded the 
highest ROC-AUC of 0.8279 among the three unsupervised models studied even though it had a false positive rate of 0.8866 when evaluated using CICIDS2017. Conversely,
all 3 unsupervised models achieved ROC-AUC scores below 0.50, implying uneffective transfer from CSE-CIC-IDS2018 to UNSW-NB15. These results, indicate that a good 
discrimination indicator does not necessarily mean a good operating point for actual use.

This type of incosistency between overall performance and successful attack detection was observed in the performance of supervised models as well.  Random Forest saw 
accuracies of 0.9188 and 0.9743 for CSE-CIC-IDS2018 and UNSW-NB15 respectively, but its attack recall were only 0.2043 and 0.0002 respectively. XGBoost saw ROC-AUC scores 
of 0.9153 and 0.7829 across the two test datasets. Despite better 
ROC-AUC scores, the recall score stood low at 0.3150 and 0.0739 implying the model's effectiveness in detecting minority attack classes.

The outcome of DNN serves as an additional proof that the configuration of a model lead to different results across datasets. The first Optuna study of DNN achieved ROC-AUC 
of 0.5954 and 0.8096 with recall score of 0.0005 and 0.0794 for the 2 test datasets. The second Optuna study showed lower scores than the first one. Despite some improvements 
in metrics, DNN still showed inconsistent performance lacking the capability to detect attacks even after optimizing the hyperparameters.

Altogether, the results suggest that NIDS evaluation across datasets should provide information on a number of supporting metrics instead of relying on a single metric. This is 
significant since accuracy can still be high even if the models detect few attacks when the datasets are highly imbalanced. Despite the usefulness of ROC-AUC for evaluating ranking 
performance, it does not indicate whether the threshold set allows for the desired number of missings and false alarms.

None of the models showed robustness in generalizing for NIDS. Hence, they are not ready for deployments. A high false-positive rate can be very demanding for security analysts in a 
real intrusion detection process, while a low recall may allow attacks to remain undetected. Hence, it is important to have a clearly determined false alarm budget, the correct calibration 
of thresholds on representative validation data, further validation using real traffic of the targe environment and an assessment of resources required for alerts investigations in order to 
deploy the model in real scenario.

To sum up, from the results, it can be concluded that predictive performance based on a dataset doesn't guarantee that the model can be used on other unseen datasets. Even though, if the model
performs well on one benchmark dataset, it doesn't mean that the model is trustworthy in generating better results from other dataset. Thus, the model selection criteria should consider classical 
aggregate performance criteria as well as attack-class recall, false-positive rate, consistent performance across datasets and the operational threshold stability. 
\section{Research Contributions}
\label{sec:conclusion-contributions}

The primary contribution is consistent directional comparison of three unsupervised anomaly-detection models across three intrusion-detection 
datasets. The experiment used a common nine-feature flow representation, source-fitted preprocessing, normal-only model training, labelled 
source-domain validation for model and threshold selection, and fixed target test partitions. This design made the nine source-to-target paths 
comparable within the recorded experiment.

A second contribution is the joint interpretation of ranking and operating-point performance. The difference between ROC-AUC and F1\_Score or FPR 
is not a technical detail: it changes the conclusion that can reasonably be drawn about transferred detector. Several models retained some ranking 
ability while producing many false alerts at their source-calibrated thresholds. Reporting all three measures prevents a favourable ROC-AUC from 
concealing poor alert quality.

A third contribution is the demonstration that no model led every transfer path. XGBoost generated the highest ROC-AUC of 0.9153 and F1\_Score of 0.4768 
for CSE-CIC-IDS2018. DNN produced the highest ROC-AUC of 0.8107, best F1\_Score of 0.1328 by XGBoost for UNSW-NB15 amongst supervised models. The 
autoencoder had the highest mean cross-dataset ROC-AUC (0.6490), while Isolation Forest had the highest mean cross-dataset F1\_Score (0.3084) amongst unsupervised models. Neither model can be called uniformly robust on that 
basis: the difference were not tested across repeated runs, and some loading paths had high FPR.

The study does not claim a new anomaly-detection architecture or overall superiority over published methods. For the CICIDS2017->CSE-CIC-IDS2018 
directions also evaluated by \cite{Verkerken2022}, the present results showed higher ROC-AUC for all the models but lower F1\_Score. It could be 
due to study variation as studies differ in feature construction, sampling, preprocessing, class composition, and threshold selection, these figures 
do not form a controlled head-to-head test. The contribution is the transparent account of what transferred, what failed, and how the answer changed 
with the metric and direction.

The Research questions can be answered as follows:

\begin{description}
       \item[R01.]
       All the supervised and unsupervised machine learning models generalised unevenly to unseen datasets. Some paths retained moderate score discrimination, whereas others had poor ROC-AUC 
       or excessive false alarms. 
       
       \item[R02.]
       Supervised Machine Learning Models outperformed Unsupervised Machine Learning Models in terms of performance for cross-dataset 
       generalization for NIDS with higher F1\_Score and ROC-AUC. 
       
       \item[R03.]
        No single model demonstrated consistently superior robustness across the dataset distribution shift in real-world network intrusion detection scenarios. XGBoost showed superior performance with F1\_Score of 0.4768 and 
        ROC-AUC score of 0.9153 while tested on CSE-CIC-IDS2018. It remained superior supervised machine learning model for UNSW-NB15 as well, however, 
        Improved Auto-Encoder showed robustness with F1\_Score of 0.3029 and ROC-AUC score of 0.7573 while tested on UNSW-NB15.
\end{description}


\section{Research Philosophy}
\label{sec:conclusion-philosophy}

The study adopted a positivist, deductive, quantitative, and comparative approach. Model type, source dataset, target dataset, anomaly scores, 
predictions, and performance measures were recorded through defined computational procedure. The conclusions are based on those observations 
rather than participant opinions.

The deductive element involved examining expectations from prior work, that performance would often deteriorate under dataset shift, that model 
behavior would differ, and that transfer direction would matter. The recorded results were consistent with these expectations descriptively. 
However, one result per model and path does not support formal inferential testing.

This approach also places a boundary on explanation. The experiment measured changes in performance but did not isolate the cause of each change. 
Differences in normal traffic, attack composition, features extraction, prevalence, and threshold suitability are plausible explanations. Their 
individual effects cannot be established from the reported results alone.

\section{Constraints and Limitations}
\label{sec:conclusion-limitations}

\begin{itemize}
    \item The three public benchmarks contained particular mixtures of traffic and attacks. They cannot and do not represent every current network 
    or future attack. Dataset coverage and realism are recognised limitations of intrusion-detection.
    \cite{Ring2019,Keyon2020}.
    \item Harmonisation required the study to use nine network-flow features that could be constructed across all three datasets. This enabled direct 
    transfer but excluded other potentially informative fields. Similar features names also do not guarantee identical measurement processes or 
    distribution across data sources.
    \item The within-dataset partitions were formed through stratified random sampling of flows. The experiment did not establish that duplicate 
    or highly similar flows were separated, nor did it evaluate a later time period as o temporal holdout. Within-dataset performance may therefore 
    be more favourable than performance under a stricter deployment-oriented split.
    \item Computational caps limited training, validation, and testing to at most 100,000 normal training flows, 50,000 validation flows, and 80,000 
    test flows per normal flows. The resulting estimates describe these sampled partitions, not every available flow.
    \item Although model fitting used normal source traffic labelled source validation data were used for model selection and threshold calibration. 
    The procedure should not be presented as entirely label-free. The same source-validation partition supported multiple development decisions, 
    potentially making validation performance optimistic. Target labels were not used for those decisions and remained reserved for final evaluation.
    \item Binary benign versus attack evaluation concealed differences among attacks categories. The experiment did not measure real-time latency, 
    resource consumption, analyst workload, or detector behavior under continuously changing production attacks.
\end{itemize} 

\section{Scope for Future Researchers}
\label{sec:conclusion-future-research}

The most immediate extension is to repeat the experiment for all the models across severals independent random seeds. Reporting path-specific 
means, dispersion, and uncertainity intervals would allow a stronger assessment of whether the observed difference are stable. A second extension 
is to use deduplicated and temporal partitions, followed by testing on traffic from operationally distinct environments. These designs would 
provide a more demanding test of whether within-dataset performance and cross-dataset ranking persist outside random flow-level splits.
Third, future studies could compare several documented common feature sets while preserving source-only statistical fitting. This would help 
distinguish information lost through nine-feature harmonisation from limitations of the anomaly-detection models. Fourth, threshold calibration 
should be evaluated explicitly. Researchers could compare source-only threshold rules and report recall at predeclared false-positive-rate budget. 
Unlabelled target adaptation may also be investigated, but it must be presented as a separate protocol. Target test labels must not determine 
model selection or thresholds.Finally, future evaluations could add precision--recall curves, average precision, attack-category analyses, and 
computational measurements. These would clarify the effects of class composition and the practical cost of operating each detector.


\section{Summary}
\label{sec:conclusion-summary}

This dissertation tested whether Supervised and Unsupervised Machine learning models trained on normal source traffic dataset CICIDS2017 generalise 
across benchmark test datasets CSE-CIC-IDS2018 and UNSW-NB15. Within-dataset performance did not reliably carry over to external datasets, and the 
direction of transfer mattered. No model showed robust generalization capability for NIDS.

The central conclusion is that cross-dataset generalisation requires more than a favourable score on one benchmark or metric. 
ROC-AUC describes transferred score discrimination. F1\_Score show the consequences of making decisions at frozen source-calibrated threshold. 
Several models retained some discrimination but generated excessive false alerts. The work provides a reproducible evaluation framework and a 
bounded account of its findings, while identifying the validation needed before claims of reliable real-world intrusion detection can be made.
